% !PRINTED 587
\weberpage{587}

If $\alpha$ is divisible by $\beta$, then $N(\alpha)$ is divisible by $N(\beta)$.
For from $\alpha=\beta\gamma$ it follows that $N(\alpha)=N(\beta)N(\gamma)$.

If $\alpha$, $\beta$, $\gamma$ are integers in $R(i)$ and $\alpha$ is divisible by $\gamma$,
then $\alpha\beta$ is also divisible by $\gamma\beta$; and if $\alpha$ and $\beta$ are divisible by $\gamma$,
then $\alpha\pm\beta$ is also divisible by $\gamma$.

Every integer is divisible by every unit.

% [source-original mathematical ellipsis, verified 2026-08-29] From
% $N(\alpha)N(\beta)=1$ with positive integral norms, both norms equal $1$;
% the print omits ``$=1$'' after the first $N(\alpha)$.
The units, however, are divisible only by units; for
if $\alpha\beta$ is a unit, then $N(\alpha)N(\beta)=1$, hence $N(\alpha)$ and
$N(\beta)=1$, i.\,e. $\alpha$ and $\beta$ are units.

Associated numbers are mutually divisible by one another,
and when two numbers are mutually divisible, they are
associated. For if $\alpha:\beta$ and $\beta:\alpha$ are both integers, then
the product of these two numbers is $=1$; hence both are units.

For the numbers of the field $R(i)$ the same
laws hold concerning the decomposition into prime factors as for the real
integers.

\gesperrt{An integer $\alpha$ of the field $R(i)$ which is not a unit
is called composite if it can be decomposed into several
integer factors, none of which is a unit.}

\gesperrt{If it cannot be so decomposed, it shall be called a
prime in the field $R(i)$.}

If $\xi=x+yi$ is a fractional number of the field $R(i)$, then
an integer $\mu=m+ni$ can be determined so that the
norm of the difference $\xi-\mu$, that is $(x-m)^2+(y-n)^2$, is smaller
than, or at least not greater than, $1/2$; for one need only
choose the integral rational numbers $m$, $n$ so that $x-m$
and $y-n$, taken absolutely, are not greater than $1/2$, and hence their squares
not greater than $1/4$.

If, then, $\alpha$ and $\alpha_1$ are two integers in $R(i)$, and $\alpha_1$ is
different from zero, one can accordingly determine the integers $\mu$
and $\alpha_2$, so that
\[
\frac{\alpha}{\alpha_1}-\mu=\frac{\alpha_2}{\alpha_1},\quad \frac{N(\alpha_2)}{N(\alpha_1)}\leqq \frac{1}{2},
\]
hence
\[
\alpha=\mu\alpha_1+\alpha_2,\quad N(\alpha_2)<N(\alpha_1).
\]
If $\alpha_2$ is not zero, one can proceed in the same way with the numbers
$\alpha_1$, $\alpha_2$ and obtains
\[
\alpha_1=\mu_1\alpha_2+\alpha_3,\quad N(\alpha_3)<N(\alpha_2),
\]
%CONT>>
